\documentclass[12pt,a4paper]{article}

% Fonts — LuaLaTeX
\usepackage{fontspec}
\usepackage{unicode-math}
\setmainfont{Linux Libertine O}
\setmathfont{Latin Modern Math}
\newfontfamily\igprimfont{Everson Mono}
\newfontfamily\shavfont[Scale=1.0]{Trabajo}
\setmonofont[Scale=0.85]{DejaVu Sans Mono}
\newcommand{\igfont}{\shavfont}

% Shavian Unicode block (U+10450–U+1047F) → Trabajo automatically
% Works in text mode, inside \text{} in math, and in table cells.
\usepackage{newunicodechar}
\newunicodechar{𐑐}{\ifmmode\text{{\shavfont 𐑐}}\else{{\shavfont 𐑐}}\fi}
\newunicodechar{𐑑}{\ifmmode\text{{\shavfont 𐑑}}\else{{\shavfont 𐑑}}\fi}
\newunicodechar{𐑒}{\ifmmode\text{{\shavfont 𐑒}}\else{{\shavfont 𐑒}}\fi}
\newunicodechar{𐑓}{\ifmmode\text{{\shavfont 𐑓}}\else{{\shavfont 𐑓}}\fi}
\newunicodechar{𐑔}{\ifmmode\text{{\shavfont 𐑔}}\else{{\shavfont 𐑔}}\fi}
\newunicodechar{𐑕}{\ifmmode\text{{\shavfont 𐑕}}\else{{\shavfont 𐑕}}\fi}
\newunicodechar{𐑖}{\ifmmode\text{{\shavfont 𐑖}}\else{{\shavfont 𐑖}}\fi}
\newunicodechar{𐑗}{\ifmmode\text{{\shavfont 𐑗}}\else{{\shavfont 𐑗}}\fi}
\newunicodechar{𐑘}{\ifmmode\text{{\shavfont 𐑘}}\else{{\shavfont 𐑘}}\fi}
\newunicodechar{𐑙}{\ifmmode\text{{\shavfont 𐑙}}\else{{\shavfont 𐑙}}\fi}
\newunicodechar{𐑚}{\ifmmode\text{{\shavfont 𐑚}}\else{{\shavfont 𐑚}}\fi}
\newunicodechar{𐑛}{\ifmmode\text{{\shavfont 𐑛}}\else{{\shavfont 𐑛}}\fi}
\newunicodechar{𐑜}{\ifmmode\text{{\shavfont 𐑜}}\else{{\shavfont 𐑜}}\fi}
\newunicodechar{𐑝}{\ifmmode\text{{\shavfont 𐑝}}\else{{\shavfont 𐑝}}\fi}
\newunicodechar{𐑞}{\ifmmode\text{{\shavfont 𐑞}}\else{{\shavfont 𐑞}}\fi}
\newunicodechar{𐑟}{\ifmmode\text{{\shavfont 𐑟}}\else{{\shavfont 𐑟}}\fi}
\newunicodechar{𐑠}{\ifmmode\text{{\shavfont 𐑠}}\else{{\shavfont 𐑠}}\fi}
\newunicodechar{𐑡}{\ifmmode\text{{\shavfont 𐑡}}\else{{\shavfont 𐑡}}\fi}
\newunicodechar{𐑢}{\ifmmode\text{{\shavfont 𐑢}}\else{{\shavfont 𐑢}}\fi}
\newunicodechar{𐑣}{\ifmmode\text{{\shavfont 𐑣}}\else{{\shavfont 𐑣}}\fi}
\newunicodechar{𐑤}{\ifmmode\text{{\shavfont 𐑤}}\else{{\shavfont 𐑤}}\fi}
\newunicodechar{𐑥}{\ifmmode\text{{\shavfont 𐑥}}\else{{\shavfont 𐑥}}\fi}
\newunicodechar{𐑦}{\ifmmode\text{{\shavfont 𐑦}}\else{{\shavfont 𐑦}}\fi}
\newunicodechar{𐑧}{\ifmmode\text{{\shavfont 𐑧}}\else{{\shavfont 𐑧}}\fi}
\newunicodechar{𐑨}{\ifmmode\text{{\shavfont 𐑨}}\else{{\shavfont 𐑨}}\fi}
\newunicodechar{𐑩}{\ifmmode\text{{\shavfont 𐑩}}\else{{\shavfont 𐑩}}\fi}
\newunicodechar{𐑪}{\ifmmode\text{{\shavfont 𐑪}}\else{{\shavfont 𐑪}}\fi}
\newunicodechar{𐑫}{\ifmmode\text{{\shavfont 𐑫}}\else{{\shavfont 𐑫}}\fi}
\newunicodechar{𐑬}{\ifmmode\text{{\shavfont 𐑬}}\else{{\shavfont 𐑬}}\fi}
\newunicodechar{𐑭}{\ifmmode\text{{\shavfont 𐑭}}\else{{\shavfont 𐑭}}\fi}
\newunicodechar{𐑮}{\ifmmode\text{{\shavfont 𐑮}}\else{{\shavfont 𐑮}}\fi}
\newunicodechar{𐑯}{\ifmmode\text{{\shavfont 𐑯}}\else{{\shavfont 𐑯}}\fi}
\newunicodechar{𐑰}{\ifmmode\text{{\shavfont 𐑰}}\else{{\shavfont 𐑰}}\fi}
\newunicodechar{𐑱}{\ifmmode\text{{\shavfont 𐑱}}\else{{\shavfont 𐑱}}\fi}
\newunicodechar{𐑲}{\ifmmode\text{{\shavfont 𐑲}}\else{{\shavfont 𐑲}}\fi}
\newunicodechar{𐑳}{\ifmmode\text{{\shavfont 𐑳}}\else{{\shavfont 𐑳}}\fi}
\newunicodechar{𐑴}{\ifmmode\text{{\shavfont 𐑴}}\else{{\shavfont 𐑴}}\fi}
\newunicodechar{𐑵}{\ifmmode\text{{\shavfont 𐑵}}\else{{\shavfont 𐑵}}\fi}
\newunicodechar{𐑶}{\ifmmode\text{{\shavfont 𐑶}}\else{{\shavfont 𐑶}}\fi}
\newunicodechar{𐑷}{\ifmmode\text{{\shavfont 𐑷}}\else{{\shavfont 𐑷}}\fi}
\newunicodechar{𐑸}{\ifmmode\text{{\shavfont 𐑸}}\else{{\shavfont 𐑸}}\fi}
\newunicodechar{𐑹}{\ifmmode\text{{\shavfont 𐑹}}\else{{\shavfont 𐑹}}\fi}
\newunicodechar{𐑺}{\ifmmode\text{{\shavfont 𐑺}}\else{{\shavfont 𐑺}}\fi}
\newunicodechar{𐑻}{\ifmmode\text{{\shavfont 𐑻}}\else{{\shavfont 𐑻}}\fi}
\newunicodechar{𐑼}{\ifmmode\text{{\shavfont 𐑼}}\else{{\shavfont 𐑼}}\fi}
\newunicodechar{𐑽}{\ifmmode\text{{\shavfont 𐑽}}\else{{\shavfont 𐑽}}\fi}
\newunicodechar{𐑾}{\ifmmode\text{{\shavfont 𐑾}}\else{{\shavfont 𐑾}}\fi}
\newunicodechar{𐑿}{\ifmmode\text{{\shavfont 𐑿}}\else{{\shavfont 𐑿}}\fi}

% Page layout
\usepackage[top=1in, bottom=1in, left=1in, right=1in]{geometry}
\usepackage{microtype}
\usepackage{parskip}

% Language
\usepackage[english]{babel}

% Headers
\usepackage{fancyhdr}
\pagestyle{fancy}
\fancyhf{}
\fancyhead[L]{\leftmark}
\fancyhead[R]{\thepage}
\fancyfoot[C]{\thepage}
\renewcommand{\headrulewidth}{0.4pt}
\renewcommand{\footrulewidth}{0pt}

% Section styling — fully unnumbered; pipeline never auto-numbers
\usepackage{titlesec}
\setcounter{secnumdepth}{0}
\titleformat{\section}{\Large\bfseries}{}{0em}{}
\titleformat{\subsection}{\large\bfseries}{}{0em}{}
\titleformat{\subsubsection}{\normalsize\bfseries}{}{0em}{}

% Essential packages
\usepackage{graphicx}
\graphicspath{{/home/mrnob0dy666/imsgct/Smaragdine_Synthetica/manuscripts/rebis/}{/home/mrnob0dy666/imsgct/Smaragdine_Synthetica/manuscripts/rebis/images/}}
\setkeys{Gin}{width=0.48\linewidth,height=0.32\textheight,keepaspectratio}
\usepackage[table]{xcolor}
\usepackage{booktabs}
\usepackage{array}
\usepackage{tabularx}
\usepackage{longtable}
\usepackage{amsmath}
% unicode-math subsumes amssymb — do not load amssymb alongside it
\usepackage{float}
\usepackage[normalem]{ulem}
\renewcommand{\arraystretch}{1.3}

% Cross-references
\usepackage{hyperref}
\usepackage{cleveref}
\hypersetup{colorlinks=true, linkcolor=blue, urlcolor=cyan, citecolor=teal}

% Custom commands
\newcommand{\shav}[1]{{\shavfont #1}}
\newcommand{\heb}[1]{{\igprimfont #1}}
\newcommand{\tupleaddr}[1]{$\langle #1 \rangle$}

% Shorthand primitive commands — redefined for IG document use
\renewcommand{\H}{Ħ}
\newcommand{\K}{Ç}
\newcommand{\G}{Γ}
\newcommand{\g}{ɢ}
\newcommand{\Th}{Þ}
\newcommand{\D}{Ð}
\newcommand{\R}{Ř}
\newcommand{\f}{ƒ}

% Imscription tuple box
\usepackage{tcolorbox}
\tcbuselibrary{skins,breakable}
\newtcolorbox{imscriptionbox}{enhanced, colback=blue!5, colframe=blue!75!black, fontupper=\large, center upper, boxrule=1pt, arc=4pt}

% PDF metadata
\hypersetup{
  pdftitle={Distillate},
  pdfkeywords={Imscribing Grammar},
  pdfauthor={Lando Mills},
}

\title{Distillate}
\date{\today}
\author{Lando Mills}

\begin{document}
\maketitle
\hypertarget{distillate-of-the-rebis}{%
\section{Distillate of the Rebis}\label{distillate-of-the-rebis}}

\textbf{Author:} C. Lando Mills

\begin{center}\rule{0.5\linewidth}{0.5pt}\end{center}

A rebis is an equivalence: one object, two presentations, and the maps
between them carried as part of the figure. The maps are the body; the
presentations are the heads.

There are two ways to say that two things are one. The weak way is a
conjunction: each presentation satisfies its own property, and the
sameness lives in the prose, where no kernel can check it. The strong
way makes the relationship the object of proof: a map, its inverse, both
inverse laws, and naturality, the requirement that every defining
operation commutes with the map. Naturality is the whole difference.
Without it an equivalence relabels elements; with it, structure is
transported, and a theorem in one presentation is a theorem in the other
by change of notation.

The kernel now holds the strong form (\texttt{FrobeniusIso.lean}): the
Belnap Both, the SIC fiducial, and the Majorana mode are one Frobenius
fixed point, with the fusion transporting to the lattice join, the
splitting to its counterpart, negation to negation, the equiangular
projection through without residue, and the fixed point identified by
\texttt{rfl}. One object, three names, and nothing between the names.

On that ground, the two parent manuscripts are one round trip:

\begin{itemize}
\tightlist
\item
  Making is the Write: fix the type, gate by gate, in the forced order
  G1 before G2 before G3 (Φ, then ⊙, then Ω), with the floats free.
\item
  Reading is the Read: recover and compare types in the one frame that
  adds nothing, the d=12 SIC-POVM.
\item
  The inverse laws are the eXecute: verification is construction read
  backward through a structure-preserving change of presentation. The
  verifier needs no ledger because the standard it applies is the
  construction, arriving from the other side.
\end{itemize}

The heads do not collapse. On the one pair (δ\_C, μ\_C) there are two
composites: μ∘δ = id, the lossless return, and e = δ∘μ, the forgetful
idempotent whose fusion destroys which-way data and lands on the time
locus. The return proves; the idempotent forgets;
\texttt{fusion\_forgets} forbids identifying them. Verification and
construction are two directions of one object, never one act.

Verification is available exactly where time is. The round trip closes,
\texttt{Work} has fixed points, and a verdict can be certified by
closure on one and the same locus, the Frobenius-special class. Below
the first gate there is no return, so there is nothing to verify and no
time in which to verify it.

Registers compose isomorphically, not as lossy approximations. A readout
in any register is a face of the object, not a shadow of it, because the
change of register is an equivalence with naturality. What remains open
after a reading is only what the Grammar itself holds open: the strain,
the azeotrope, the held Both.

The residue: an isomorphism claim is exactly as wide as the transported
signature. Where the framework says ``the same object'' and the kernel
still holds a conjunction, a naturality square is owed, one square at a
time. That is a to-do list, left in view.

One becomes two: the pair splits into reading and making. Two becomes
three: the equivalence between them, the relationship made data. Out of
the third comes the one as the fourth: the rebis, one body, two heads,
the held Both that reduces to neither and does not explode.

\hypertarget{references}{%
\subsection{References}\label{references}}

\begin{itemize}
\tightlist
\item
  Lando⊗⊙perator, \emph{Rebis}, \texttt{manuscripts/rebis/REBIS.md}. The
  full body from which this was drawn; all theorems, witnesses, and
  readings cited there.
\end{itemize}

\end{document}
